% =====================================================================
%  Test Vector Generation for Fault Detection using GNNs
%  BTech Mini Project Report
%  Authors : Prasham Shah     (202301252)
%            Dhruvi Jobanputra (202301401)
%  School  : Dhirubhai Ambani University
% =====================================================================
\documentclass[11pt,a4paper]{article}

% ---------- packages -------------------------------------------------
\usepackage[utf8]{inputenc}
\usepackage[T1]{fontenc}
\usepackage[margin=1in]{geometry}
\usepackage{graphicx}
\usepackage{amsmath,amssymb}
\usepackage{booktabs}
\usepackage{caption}
\usepackage{microtype}
\usepackage{enumitem}
\usepackage{xcolor}
\usepackage[colorlinks=true,linkcolor=black,citecolor=black,urlcolor=blue!60!black]{hyperref}

% ---------- figure path ----------------------------------------------
\graphicspath{{visualization/figures/}{report-result-images/}}

% ---------- helpers --------------------------------------------------
\newcommand{\R}{\mathbb{R}}
\newcommand{\code}[1]{\texttt{#1}}
\newcommand{\bd}[1]{\textbf{#1}}
\captionsetup{font=small,labelfont=bf,labelsep=period}

% ---------- title ----------------------------------------------------
\title{\bfseries Test Vector Generation for Fault Detection\\
using Graph Neural Networks}

\author{Prasham Shah (202301252) \quad Dhruvi Jobanputra (202301401)\\[4pt]
\small BTech Mini Project\\
\small School of Engineering and Applied Science\\
\small Dhirubhai Ambani University}

\date{\today}

% =====================================================================
\begin{document}
\maketitle

% =====================================================================
\begin{abstract}
\noindent
Automatic Test Pattern Generation (ATPG) is the problem of finding an input
assignment that exposes a specific manufacturing fault at a circuit's
primary outputs. Classical ATPG algorithms (D-Algorithm, PODEM, FAN) solve
it via tree search with backtracking, which scales poorly on large
reconvergent circuits. This report documents a from-scratch reproduction
of \emph{DETECTive} (Petrolo, Medya, Graziano, Pal --- GLSVLSI '24), a
deep-learning model that replaces the search with a single forward pass: a
graph neural network (GAT + GCN stack) encodes the faulty netlist, two
LSTMs encode the fault's activation and propagation paths, and a final MLP
predicts each primary input's Boolean value. We re-implement the entire
pipeline in PyTorch and PyTorch Geometric, regenerate the dataset of
$\sim$15{,}000 fault-pattern pairs from 2{,}000 random 4-input circuits, and
evaluate against the paper's claims on bit accuracy, fault depth, and
reconvergent fanout. Our model reproduces the paper's headline numbers in
the per-bucket regimes the paper reports against
(shallow-fault accuracy 0.88 vs. paper's $\sim$0.82, deep-fault accuracy
0.80 vs. paper's $\sim$0.66, low-reconvergence accuracy 0.89 vs. paper's
$\sim$0.88). 50.6\,\% of validation samples are predicted with perfect
bit accuracy. The report is organised around 17 architectural block
diagrams that walk a reader through the full pipeline, from a faulted
Verilog netlist to the predicted primary-input vector.
\end{abstract}

\tableofcontents
\bigskip

% =====================================================================
\section{Introduction}
\label{sec:intro}

% --- 1.1 motivation -------------------------------------------------
\subsection{The ATPG problem}

A digital integrated circuit is manufactured from millions of transistors,
and any one of those transistors can fail. The dominant model used in
production test today is the \emph{stuck-at} fault model: a wire is
permanently held at logic 0 (\bd{stuck-at-0}) or logic 1
(\bd{stuck-at-1}), regardless of what should drive it. To detect such a
fault, the test engineer must find a primary-input vector that
\emph{(a)}~drives the faulty wire to the opposite of its stuck value
(activation), and \emph{(b)}~propagates the resulting logical difference
through downstream logic to a primary output where it can be observed
(propagation). A circuit with $n$ primary inputs admits $2^{n}$
candidate vectors; for the largest ICs in production, $n$ is in the
hundreds of thousands and exhaustive enumeration is infeasible.

Figure~\ref{fig:atpg} states the problem in one diagram: given a
\emph{faulty} netlist (one wire annotated with a stuck-at value), find an
input vector $\hat{p}$ such that simulating the circuit under that vector
produces a logic value at some PO that differs from the fault-free
simulation.

\begin{figure}[htbp]
  \centering
  \includegraphics[width=0.95\linewidth]{fig01_atpg_problem}
  \caption{The ATPG problem as a constraint-satisfaction task. The
  classical solvers (D-Algorithm, PODEM, FAN) tackle it with depth-first
  search and backtracking; DETECTive learns a direct mapping from
  faulted netlist to test pattern.}
  \label{fig:atpg}
\end{figure}

% --- 1.2 classical ATPG --------------------------------------------
\subsection{Classical ATPG and its limits}

The three milestone algorithms are \bd{D-Algorithm} (Roth, 1966),
\bd{PODEM} (Goel, 1981), and \bd{FAN} (Fujiwara \& Shimono, 1983). All
three pose ATPG as tree search:
\begin{itemize}[leftmargin=*]
  \item assign a value to one PI;
  \item simulate forward / justify backward;
  \item if a contradiction is found, undo the most recent decision and
        try the other branch (\emph{backtracking}).
\end{itemize}
On combinational circuits with reconvergent fanout, the search tree
grows exponentially and the worst-case complexity is NP-complete. PODEM
and FAN reduce the average constant factor (only PI variables enter the
decision stack; FAN further uses headlines to limit backtracking depth),
but the asymptotic behaviour is unchanged.

% --- 1.3 ML-based ATPG ----------------------------------------------
\subsection{ML-based ATPG and DETECTive's contribution}

DETECTive's idea is to side-step the search entirely. It treats the
faulted circuit as a labelled, directed graph and learns a function
\[
  f_\theta : (G, \text{fault}) \;\mapsto\; \hat{p} \in [0,1]^{n}
\]
where $G$ is the graph encoding of the netlist and $\hat{p}$ is a
probability vector over the primary inputs. The model is trained
end-to-end with binary cross-entropy against ground-truth patterns
generated by exhaustive simulation on small circuits ($n \!\le\! 4$).
Because $f_\theta$ is a single forward pass through three small neural
modules, inference is constant in the size of the search space and
embarrassingly parallel across faults.

The remainder of this report unpacks the model layer by layer. Each
component is illustrated with a self-contained block diagram (Figures
\ref{fig:pipeline}--\ref{fig:walk}), and the data shape is annotated at
every interface so a reader can trace a single fault from a Verilog
string to a Boolean PI vector.

% =====================================================================
\section{Background and Related Work}
\label{sec:bg}

\paragraph{Classical algorithms.} D-Algorithm, PODEM and FAN are the
three milestones of search-based ATPG. We give each a dedicated
subsection in Section~\ref{sec:classical}, since we implement all
three of them in this project and benchmark them against DETECTive.

\paragraph{Graph Neural Networks.} A GNN learns one embedding vector per
graph node by repeatedly aggregating each node's feature with the
features of its neighbours. The two flavours used by DETECTive are
\bd{GAT} (Graph Attention Networks --- Veli\v{c}kovi\'c et al., 2018),
which weights each neighbour by a learned attention coefficient, and
\bd{GCN} (Graph Convolutional Networks --- Kipf \& Welling, 2017),
which weights each neighbour by a normalised adjacency. Both are
implemented in PyTorch Geometric and we use the library's reference
implementations directly.

\paragraph{LSTM.} A Long Short-Term Memory cell (Hochreiter \&
Schmidhuber, 1997) processes a variable-length sequence with a hidden
state $h_t$ and a cell state $c_t$, gated by three sigmoid gates
(forget, input, output). DETECTive uses LSTMs to compress an
arbitrary-length \emph{path} (a sequence of node embeddings along a
PI$\to$fault or fault$\to$PO chain) into a single fixed-size vector.

% =====================================================================
\section{Classical ATPG Algorithms}
\label{sec:classical}

Three algorithms dominate the classical ATPG literature: the
\bd{D-Algorithm} (Roth, 1966), \bd{PODEM} (Goel, 1981) and \bd{FAN}
(Fujiwara \& Shimono, 1983). Each was a strict improvement over its
predecessor --- algorithmically, in the size of the decision space,
and in the number of backtracks needed on real circuits. Together they
define the search-based family that DETECTive is built to compete
with. We implement all three of them and benchmark them on ISCAS-85
designs alongside DETECTive (Section~\ref{sec:eval}). The remainder of this section walks through
how each algorithm works and what made the next one a worthwhile
improvement.

% --------------------------------------------------------------------
\subsection{D-Algorithm (Roth, 1966): the first systematic ATPG}

The D-Algorithm was the first method that posed ATPG as a well-defined
constraint-propagation problem rather than ad-hoc simulation. Its
central idea is a five-valued logic in which the \emph{fault effect}
travels through the circuit alongside ordinary Boolean values:
\[
  \{0,\; 1,\; X,\; D,\; \bar{D}\}
\]
where $X$ is ``unassigned / don't care,''
$D$ is ``1 in the fault-free circuit but 0 in the faulty,'' and
$\bar D$ is the opposite. A test pattern exists when at least one $D$
or $\bar D$ symbol reaches a primary output; the difference between
the fault-free and faulty circuits is then observable.

\paragraph{Three phases.} Given a fault, the algorithm runs three
logical phases:
\begin{enumerate}[leftmargin=*,itemsep=2pt]
  \item \bd{Fault excitation.} Assign $D$ (or $\bar D$, depending on
        the stuck-at value) to the faulty line. Without this, the
        faulty circuit is identical to the fault-free one --- the
        fault is dormant.
  \item \bd{D-drive (forward propagation).} Push the fault effect
        forward toward a primary output. The set of unprocessed gates
        whose inputs include $D/\bar D$ but whose output is still $X$
        is the \emph{D-frontier}. To extend the frontier through a
        gate, the algorithm assigns the gate's other inputs the
        non-controlling value (e.g., 1 for AND, 0 for OR), so the
        $D$ on the controlling input survives the gate.
  \item \bd{Line justification (backward).} Every line forced during
        D-drive must be traced back to a consistent PI assignment.
        Whenever an internal line is forced to a value, the algorithm
        recurses on the gate that drives it. If a contradiction
        appears (a line forced to both 0 and 1), the algorithm
        backtracks to the most recent decision and tries the
        alternative.
\end{enumerate}

\paragraph{Decision space.} D-Algorithm permits decisions on
\emph{any} internal line, not just primary inputs. The worst-case
decision tree therefore has up to $2^{N}$ leaves, where $N$ is the
number of nodes (gates plus inputs). Reconvergent fanout is
particularly painful: a PI feeds two paths that meet again
downstream, and the conflict between paths is only discovered after
both branches have been explored.

\paragraph{Limitations.} The combination of internal-line decisions
and full backward justification means D-Algorithm scales poorly with
size and reconvergence. On ISCAS circuits like c2670 or c7552, more
than 99\,\% of runtime can be spent on backtracks. This motivated
PODEM.

% --------------------------------------------------------------------
\subsection{PODEM (Goel, 1981): decisions only on primary inputs}

Goel's insight was that every internal line value is, in the end, a
\emph{consequence} of the PI assignment via plain forward simulation.
If decisions are restricted to PIs, then internal contradictions are
always caught by simulation alone --- no separate line-justification
phase is needed.

PODEM (``\textbf{P}ath-\textbf{O}riented \textbf{D}ecision
\textbf{M}aking'') replaces the three-phase D-Algorithm with a single
loop:
\begin{quote}\small
\ttfamily
while fault\_effect not at a PO and not exhausted:\\
\hspace*{2em}objective $\leftarrow$ pick a line that needs a value to make progress\\
\hspace*{2em}\hspace*{2em}// either to activate the fault or to extend the D-frontier\\
\hspace*{2em}(pi, val) $\leftarrow$ backtrace(objective)~~~// walk back to a PI\\
\hspace*{2em}assign pi := val\\
\hspace*{2em}forward-simulate the entire circuit\\
\hspace*{2em}if D reached a PO:\\
\hspace*{2em}\hspace*{2em}return current PI assignment\\
\hspace*{2em}if D-frontier became empty (fault died):\\
\hspace*{2em}\hspace*{2em}backtrack: flip the most recent PI and re-simulate
\end{quote}

\paragraph{Why this is faster.} The decision tree has at most $2^{n}$
leaves where $n$ is the number of \emph{primary inputs}, not the
number of nodes. For ISCAS c7552 ($n \approx 207$ PIs and $\sim$3{,}500
gates), this is a $2^{3300}$-fold reduction in worst-case search
space. In practice the speedup is around an order of magnitude on
typical designs.

\paragraph{Backtrace heuristic.} When an objective could be satisfied
by several different PIs, PODEM picks the ``cheapest'' one according
to a controllability cost function. This is purely a heuristic ---
different costs change runtime but not correctness. ATALANTA, the
open-source ATPG tool we use as a runtime baseline, is a
PODEM-derived implementation with carefully tuned cost functions.

\paragraph{Limitations.} PODEM still re-simulates the whole circuit
after every PI assignment, and it backtraces objectives one at a time.
On heavily reconvergent circuits, the same objective can fire
repeatedly and each retry is $O(\text{circuit size})$. This motivated
FAN.

% --------------------------------------------------------------------
\subsection{FAN (Fujiwara \& Shimono, 1983): smarter backtrace}

FAN (``\textbf{Fan}out-oriented test generation'') adds two
optimisations on top of PODEM that, in combination, reduce average
backtrack counts by another factor of 5--10$\times$ on ISCAS designs.

\paragraph{1. Multiple backtrace.} Rather than backtracing one
objective at a time, FAN maintains a \emph{set} of current objectives
and propagates them backward together. When two objectives reach the
same fanout stem with conflicting requirements (one wants 0, the other
wants 1), FAN detects the conflict immediately --- before any PI is
assigned --- and backtracks at once. PODEM, by contrast, would have
made one or both PI decisions and only re-discovered the conflict
later by simulation.

\paragraph{2. Headlines.} A \emph{headline} is the output of a
\emph{fanout-free} sub-circuit: a sub-region whose only outgoing fanout
is the headline itself. Inside a fanout-free region, every line value
is uniquely determined by the headline value, so there is no need to
backtrack inside it. FAN exploits this by treating headlines as
``virtual PIs'' during backtrace:
\begin{itemize}[leftmargin=*,itemsep=2pt]
  \item Backtrace stops as soon as it reaches a headline.
  \item The headline value is fixed and pushed onto the decision stack.
  \item Justification through the headline's fanin sub-circuit is
        \emph{deferred} until a $D$ has actually reached a PO.
\end{itemize}
By the time deferred justification runs, many initial decisions have
already been confirmed; the final pass is usually a one-shot
simulation rather than another backtracking phase.

\paragraph{Net effect.} FAN was the de-facto reference algorithm for
two decades and is the backbone of most academic and commercial ATPG
suites. The runtime gap between FAN and PODEM is the main reason
ATALANTA (the tool we benchmark against in
Section~\ref{sec:eval}) is so much faster than a textbook PODEM.

% --------------------------------------------------------------------
\subsection{Where DETECTive fits}
\label{sec:classical-vs-detective}

Table~\ref{tab:atpg-progression} summarises the progression. Within
the classical search-based family, each row is a strict improvement
over the row above. DETECTive sits outside this family: it replaces
the search itself with a learned mapping from faulted netlist to
candidate test pattern, so the notion of ``backtracks'' simply does
not apply.

\begin{table}[htbp]
\centering
\caption{Progression of classical ATPG algorithms (1966--1983) and
where DETECTive fits relative to them. The ``backtracks'' column gives
the relative reduction observed on ISCAS-85 benchmarks.}
\label{tab:atpg-progression}
\small
\begin{tabular}{@{}lllll@{}}
\toprule
\bd{Algorithm} & \bd{Year} & \bd{Decisions on}      & \bd{Key innovation}             & \bd{Backtracks (relative)} \\
\midrule
D-Algorithm    & 1966      & all internal lines      & 5-valued logic ($D$, $\bar D$)  & baseline (very high)       \\
PODEM          & 1981      & primary inputs only     & path-oriented decisions         & $\sim$10$\times$ fewer     \\
FAN            & 1983      & PIs + headlines         & multi-objective backtrace,      & $\sim$5--10$\times$ fewer  \\
               &           &                         & fanout-free regions             & than PODEM                 \\
\bd{DETECTive} & 2024      & --- (forward pass)      & GNN + LSTM learned mapping      & \bd{0}                     \\
\bottomrule
\end{tabular}
\end{table}

In our experiments, all four are run side-by-side on c17, c432,
c499 and c880. The classical algorithms produce a guaranteed-correct
test pattern (or proof of redundancy) at the cost of seconds-to-minutes
of search; DETECTive produces a candidate pattern in
2--5\,ms per fault but with the per-bucket bit-accuracy reported in
Section~\ref{sec:eval}. The two approaches are therefore complementary:
DETECTive is a fast first-pass that can be cross-checked, or used as a
warm-start, by ATALANTA-style search.

% =====================================================================
\section{An Illustrative Example}
\label{sec:example}

To anchor the rest of the report, we will track one small fault all the
way through the pipeline. Figure~\ref{fig:example} shows a 4-input
combinational netlist that computes $y = \overline{ab \wedge (c+d)}$,
with the gate \code{AND1} held \emph{stuck-at-0}. The same circuit and
fault appear in every subsequent block diagram, with shape annotations
showing how its representation changes at each stage.

\begin{figure}[htbp]
  \centering
  \includegraphics[width=0.92\linewidth]{fig02_example_circuit}
  \caption{Running example. A 4-input circuit (\code{a, b, c, d}) with a
  stuck-at-0 fault on \code{AND1}. To detect it we must simultaneously
  drive \code{AND1}'s output high (so the fault makes a difference) and
  propagate that difference to the PO \code{y}. The orange edges mark
  the activation paths; the purple edge marks the propagation path.}
  \label{fig:example}
\end{figure}

\paragraph{What a correct test pattern looks like.}
For \code{AND1} stuck-at-0 to be detectable at \code{y}, we need
$ab=1$ (activation) and $c \!+\! d = 1$ (propagation through
\code{AND2}). Two of the four PIs are free, so there are six valid
patterns: $\{1101, 1110, 1111, \ldots\}$. Brute-force on a 4-input
circuit (16 vectors) enumerates them in microseconds, which is exactly
how DETECTive's training labels are generated.

% =====================================================================
\section{Pipeline Overview}
\label{sec:pipeline}

The full pipeline is summarised in Figure~\ref{fig:pipeline}. The five
high-level stages are:
\begin{enumerate}[leftmargin=*,itemsep=2pt]
  \item \bd{Graph construction.} Parse the Verilog netlist; build a
        directed graph whose nodes are gates / PIs and whose edges
        point from inputs to gates. Annotate the faulty node.
  \item \bd{GNN encoding.} A GAT layer followed by two GCN layers
        produces a 32-dimensional embedding per node. The embedding
        captures both the local gate identity and the surrounding
        structure --- crucial for ATPG, where a wire's
        \emph{controllability} depends on its fanin cone and its
        \emph{observability} depends on its fanout cone.
  \item \bd{Path extraction.} For each PI, a BFS+capped-DFS walk yields
        the set of activation paths PI$\to$fault. From the fault, an
        analogous walk yields propagation paths fault$\to$PO. Each
        path becomes a sequence of node embeddings.
  \item \bd{LSTM compression.} The Activator LSTM compresses each
        activation path into a single 32-dim vector; the Propagator
        LSTM does the same for propagation paths. Each LSTM is followed
        by a small MLP that emits a scalar ``activation/propagation
        likelihood'' for that path.
  \item \bd{Cone aggregation and PI prediction.} Up to $P\!=\!10$
        activation paths and $P\!=\!10$ propagation paths are
        zero-padded and concatenated into a 660-dim ``cone embedding''.
        A two-layer MLP maps the cone to a single sigmoid probability,
        the predicted Boolean value of the PI under consideration.
\end{enumerate}

\begin{figure}[htbp]
  \centering
  \includegraphics[width=\linewidth]{fig03_pipeline_overview}
  \caption{End-to-end DETECTive pipeline. Tensor shapes are annotated at
  every interface. $N$ is the number of nodes in the circuit; $L_{a}$
  and $L_{p}$ are activation and propagation path lengths
  (variable). The final output is one Boolean per primary input.}
  \label{fig:pipeline}
\end{figure}

% =====================================================================
\section{Stage 1 --- From Verilog Netlist to Graph}
\label{sec:graph}

The first step turns a Verilog string into a PyTorch Geometric
\code{Data} object. A custom Verilog parser handles this step; the
result is a directed graph in which:
\begin{itemize}[leftmargin=*,itemsep=2pt]
  \item every \emph{primary input} is a node of type \code{INPUT};
  \item every \emph{logic gate} is a node of type
        \code{AND / NAND / OR / NOR / XOR / XNOR / NOT / BUF};
  \item every \emph{wire} from a gate's input to that gate is a directed
        edge.
\end{itemize}
Figure~\ref{fig:c2g} shows the conversion on the running example.

\begin{figure}[htbp]
  \centering
  \includegraphics[width=0.95\linewidth]{fig04_circuit_to_graph}
  \caption{Verilog to graph. Each gate becomes one node; the gate type
  is one-hot-encoded. Edges are oriented input $\rightarrow$ gate. The
  faulty gate is highlighted in red.}
  \label{fig:c2g}
\end{figure}

\subsection{Node features}

Each node carries an 11-dimensional feature vector
\[
  x_v \;=\; \underbrace{\mathrm{onehot}(\text{type})}_{\text{9 dims}}
            \;\Vert\;
            \underbrace{\mathbb{1}[v \text{ is faulty}]}_{\text{1 dim}}
            \;\Vert\;
            \underbrace{\log(1+\mathrm{fanout}(v))}_{\text{1 dim}}
\]
totalling $\code{FEATURE\_DIM} = 11$ dimensions. The
\code{is\_faulty} flag is the only place where
the fault location enters the model --- everything downstream is purely
structural. The \code{log1p} of fanout was chosen by the paper as a
soft proxy for ``how observable'' a node is. Figure~\ref{fig:feat}
shows the feature matrix $X \in \R^{N\times 11}$ for the running
example.

\begin{figure}[htbp]
  \centering
  \includegraphics[width=0.95\linewidth]{fig05_node_features}
  \caption{Node feature matrix $X \in \R^{N\times 11}$. Each row is one
  node; columns 1--9 are the gate-type one-hot, column 10 is the
  is\_faulty flag (1 for \code{AND1}), column 11 is the log-fanout.}
  \label{fig:feat}
\end{figure}

\subsection{Adjacency representation}

PyTorch Geometric stores edges as a 2-row tensor
\code{edge\_index} $\in \R^{2 \times E}$ in COO format (the first row
holds source indices, the second holds destinations). The full
adjacency matrix $A \in \{0,1\}^{N \times N}$ is never materialised;
the GNN layers consume \code{edge\_index} directly through
sparse-tensor primitives. Figure~\ref{fig:adj} shows both
representations side-by-side for the running example.

\begin{figure}[htbp]
  \centering
  \includegraphics[width=\linewidth]{fig06_adjacency}
  \caption{Two equivalent views of the graph topology: the COO
  edge-list (left, what the code stores) and the dense adjacency matrix
  $A$ (right, useful as a mental model). For an 8-node circuit the COO
  form is $2\times 7$; the dense form is $8\times 8$ with 7 ones.}
  \label{fig:adj}
\end{figure}

% =====================================================================
\section{Stage 2 --- The Graph Neural Network Encoder}
\label{sec:gnn}

\subsection{Why a GNN, not a CNN or RNN?}

A logic netlist is not a grid (so 2D convolutions don't apply) and not
a sequence (so a vanilla RNN would lose the topology). The natural
inductive bias is ``a node's behaviour depends on the gates it can
reach,'' which is exactly the inductive bias of a message-passing GNN.

DETECTive uses a 3-layer stack: one \bd{GAT} layer at the bottom,
followed by two \bd{GCN} layers (Figure~\ref{fig:gnnov}). The choice
mirrors the paper's Section~4.2: GAT first, because attention can learn
to up-weight the few neighbours that drive a gate's output (e.g., the
two inputs of an AND, ignoring the rest of the cone); then GCN, because
once the local feature is sharpened, simple normalised aggregation is
enough to spread the signal further.

\begin{figure}[htbp]
  \centering
  \includegraphics[width=0.95\linewidth]{fig07_gnn_overview}
  \caption{The three GNN layers. Input: $X \in \R^{N \times 11}$.
  Output: $H \in \R^{N \times 32}$, one 32-dim embedding per node.
  The hidden dimension is fixed at 32 throughout the network.}
  \label{fig:gnnov}
\end{figure}

\subsection{GAT layer math}

For a node $v$ with neighbours $\mathcal{N}(v)$, the attention layer
computes
\begin{align}
  e_{vu}    &= \mathrm{LeakyReLU}\!\big(\mathbf{a}^{\top}\,
              [W h_v \,\Vert\, W h_u]\big),
              \quad u \in \mathcal{N}(v)\cup\{v\} \\[2pt]
  \alpha_{vu} &= \frac{\exp(e_{vu})}{\sum_{k\in \mathcal{N}(v)\cup\{v\}}
              \exp(e_{vk})} \\[2pt]
  h'_v       &= \sigma\!\Big(\sum_{u\in \mathcal{N}(v)\cup\{v\}}
              \alpha_{vu}\,W h_u\Big).
\end{align}
$W \in \R^{32\times 11}$ is the per-layer weight; $\mathbf{a}\in \R^{64}$
is the attention vector; $\sigma$ is ReLU. Each neighbour $u$ ends up
contributing a weight $\alpha_{vu}$ that the model learned, summing to
1 over $v$'s neighbourhood. Figure~\ref{fig:gat} shows the data flow
through one such block.

\begin{figure}[htbp]
  \centering
  \includegraphics[width=\linewidth]{fig08_gat_math}
  \caption{GAT layer arithmetic on the running example. The attention
  scores $\alpha$ are computed from the concatenated query-key
  features, softmaxed over each node's neighbourhood, and used as
  weights in a sum. The output is one 32-dim embedding per node.}
  \label{fig:gat}
\end{figure}

\subsection{GCN layer math}

The two GCN layers are simpler: they replace the learned attention
weights with the symmetrically normalised adjacency $\tilde{A}$:
\[
  \tilde{A} \;=\; D^{-1/2}\,(A + I)\,D^{-1/2}
  \quad\Longrightarrow\quad
  H^{(k+1)} \;=\; \sigma\!\big(\tilde{A}\,H^{(k)}\,W^{(k)}\big).
\]
$D$ is the degree matrix of $A+I$ (self-loops added) and
$W^{(k)} \in \R^{32\times 32}$. Figure~\ref{fig:gcn} unpacks the matrix
multiplications block by block.

\begin{figure}[htbp]
  \centering
  \includegraphics[width=\linewidth]{fig09_gcn_math}
  \caption{One GCN layer factorises into three small matrix products:
  $H^{(k)} W^{(k)}$ (per-node linear projection), then
  $\tilde{A}\,(H^{(k)} W^{(k)})$ (one round of message passing), then a
  ReLU non-linearity (only the first GCN; the last layer leaves logits
  unscaled).}
  \label{fig:gcn}
\end{figure}

\subsection{Receptive field}

After three layers, each node's embedding incorporates information
from all nodes within a graph distance of three. For our 4-input
training circuits (max depth $\sim$4) this typically covers the entire
fanin and fanout cones, which is precisely the structural context an
ATPG decision needs. Figure~\ref{fig:rf} visualises the receptive
field of \code{AND1} expanding by one hop at each layer.

\begin{figure}[htbp]
  \centering
  \includegraphics[width=\linewidth]{fig10_receptive_field}
  \caption{Receptive-field growth across the three GNN layers. After
  layer~1, each node sees its immediate predecessors. After layer~2,
  it sees grand-parents. After layer~3, the embedding at \code{AND1}
  has already mixed information from all four PIs and the downstream
  \code{NOT/AND2} gates.}
  \label{fig:rf}
\end{figure}

% =====================================================================
\section{Stage 3 --- Path Extraction}
\label{sec:paths}

The GNN gives us a 32-dim vector per node, but ATPG is fundamentally a
\emph{path-based} problem: an input affects the fault only if there is
a path from the PI to the faulty gate (\bd{activation path}), and the
fault affects an output only if there is a path from the fault to a PO
(\bd{propagation path}). DETECTive therefore enumerates these paths
explicitly.

\paragraph{Algorithm.} A backward BFS from the fault enumerates
candidate fanin nodes; for each PI in the fanin we run a depth-capped
DFS from the PI to the fault, listing all simple paths up to a fixed
maximum length of 50 nodes (the cap prevents exponential blow-up on
heavily reconvergent topologies). The propagation paths are obtained
symmetrically, with a forward BFS from the fault and a depth-capped
DFS to each PO.

\paragraph{Output format.} Each path is a list of node indices that we
later translate into a sequence of GNN embeddings,
$h_{p_1}, h_{p_2}, \ldots, h_{p_L} \in \R^{32}$. The final tensor that
goes into the LSTM has shape $[1, L, 32]$ where $L$ is the path length.
Figure~\ref{fig:paths} shows the activation and propagation paths
extracted for the running example.

\begin{figure}[htbp]
  \centering
  \includegraphics[width=\linewidth]{fig11_path_extraction}
  \caption{The two activation paths $a\!\to\!\code{AND1}$ and
  $b\!\to\!\code{AND1}$ and the propagation path
  $\code{AND1}\!\to\!\code{AND2}\!\to\!y$. Each path becomes a tensor
  of node embeddings consumed by the corresponding LSTM.}
  \label{fig:paths}
\end{figure}

% =====================================================================
\section{Stage 4 --- LSTM Path Encoders}
\label{sec:lstm}

\subsection{LSTM cell internals}

The Activator and Propagator both use one-layer LSTMs with hidden size
$h\!=\!32$. To keep the report self-contained, Figure~\ref{fig:lstm}
walks through the gate equations (using the shorthand
$v_t = [h_{t-1};\,x_t] \in \R^{64}$ for the concatenated input):
\begin{align}
  f_t &= \sigma(W_f v_t + b_f) \\
  i_t &= \sigma(W_i v_t + b_i) \\
  \tilde{c}_t &= \tanh(W_c v_t + b_c) \\
  o_t &= \sigma(W_o v_t + b_o) \\
  c_t &= f_t \odot c_{t-1} + i_t \odot \tilde{c}_t \\
  h_t &= o_t \odot \tanh(c_t).
\end{align}

\begin{figure}[htbp]
  \centering
  \includegraphics[width=\linewidth]{fig12_lstm_cell}
  \caption{One LSTM cell, unrolled. Each gate is one $\sigma$/$\tanh$
  applied to a $64\!\to\!32$ linear layer; the four learnable matrices
  $(W_f, W_i, W_c, W_o)$ live inside the PyTorch \code{nn.LSTM}
  module.}
  \label{fig:lstm}
\end{figure}

\subsection{Activator and Propagator}

The Activator (one for the model, applied to every activation path)
takes the path tensor $\in \R^{1\times L_a\times 32}$, runs the LSTM,
keeps the final hidden state $h_{L_a} \in \R^{32}$, and concatenates a
single scalar carrying the stuck-at value (0.0 for SA0, 1.0 for SA1)
before passing through a 33$\to$16$\to$1 MLP. The MLP's scalar output
is interpreted as ``how likely it is that this PI assignment activates
the fault.''

The Propagator is identical but operates on propagation paths and does
not need the fault-type scalar (propagation behaviour is independent of
the stuck value). Figure~\ref{fig:plstms} shows both LSTMs running on
the running example.

\begin{figure}[htbp]
  \centering
  \includegraphics[width=\linewidth]{fig13_path_lstms}
  \caption{Each path is folded by an LSTM into one 32-dim encoding.
  The Activator sees fault-type as a side-channel; the Propagator does
  not. Each path also produces a scalar ``likelihood'' that is later
  packed into the cone.}
  \label{fig:plstms}
\end{figure}

% =====================================================================
\section{Stage 5 --- Cone Embedding and PI Prediction}
\label{sec:cone}

\subsection{Building the cone embedding}

Different PIs have different numbers of relevant paths. To keep the
input to the final MLP a fixed size, DETECTive uses a slot-based
scheme:
\begin{itemize}[leftmargin=*,itemsep=2pt]
  \item Reserve $P\!=\!10$ slots each for activation and propagation
        paths.
  \item Each slot stores 33 numbers: the 32-dim path encoding +
        the 1-dim scalar likelihood.
  \item If a PI has fewer than 10 paths, the empty slots are filled
        with zeros.
  \item If a PI has more than 10 paths, the top 10 are kept (in code
        order; the paper does not specify a ranking).
\end{itemize}
Concatenating both streams yields a fixed
$\code{cone\_dim} = P\,(h+1)\,\cdot 2 = 10\cdot 33\cdot 2 = \bd{660}$
dimensional vector --- the \emph{cone embedding}.
Figure~\ref{fig:cone} shows the slot layout.

\begin{figure}[htbp]
  \centering
  \includegraphics[width=\linewidth]{fig14_cone_embedding}
  \caption{The cone embedding concatenates 10 activation slots and 10
  propagation slots, each 33 wide. Empty slots are zero-padded so the
  output dim is the same regardless of the number of real paths.}
  \label{fig:cone}
\end{figure}

\subsection{The MLP head}

The final head is a vanilla 2-layer MLP:
\[
  \hat{p}_i \;=\;
  \sigma\!\Big(W_2\,\mathrm{ReLU}(W_1\,c_i + b_1) + b_2\Big),
  \qquad W_1 \in \R^{64\times 660},\;\; W_2 \in \R^{1\times 64}.
\]
$c_i$ is the cone embedding of PI~$i$; $\hat{p}_i \in [0,1]$ is the
predicted probability that PI~$i$ should be set to 1 in the test
pattern. Figure~\ref{fig:mlp} shows the dimension reductions.

\begin{figure}[htbp]
  \centering
  \includegraphics[width=0.95\linewidth]{fig15_mlp_head}
  \caption{The InputPredictor MLP: $660 \to 64 \to 1$. The output is
  thresholded at 0.5 to obtain the Boolean test bit for that primary
  input.}
  \label{fig:mlp}
\end{figure}

% =====================================================================
\section{End-to-End Dataflow}
\label{sec:dataflow}

Figure~\ref{fig:flow} stitches every stage we have seen into a single
diagram, showing the tensor shape at every junction. Note in
particular:
\begin{itemize}[leftmargin=*,itemsep=2pt]
  \item the GNN is run \emph{once} per fault, producing $H \in \R^{N
        \times 32}$;
  \item the Propagator LSTM is also run once per fault (propagation
        paths do not depend on the PI under consideration);
  \item the Activator LSTM is run \emph{per PI}, on the activation
        paths starting at that PI;
  \item the InputPredictor is then run per PI to emit one Boolean.
\end{itemize}

\begin{figure}[htbp]
  \centering
  \includegraphics[width=\linewidth]{fig16_dataflow}
  \caption{Full dataflow with shapes. $N$ is the number of nodes;
  $L_a, L_p$ are activation/propagation path lengths; $K_a, K_p$ are
  the number of paths (capped at 10).}
  \label{fig:flow}
\end{figure}

\subsection{A worked example, end-to-end}

Figure~\ref{fig:walk} traces the running example through the entire
pipeline, with the actual numerical shapes filled in:
$N\!=\!8$ nodes, $K_a\!=\!2$ activation paths each of length 2,
$K_p\!=\!1$ propagation path of length 3, four PIs, and a final output
$\hat{p} \in [0,1]^{4}$.

\begin{figure}[htbp]
  \centering
  \includegraphics[width=\linewidth]{fig17_walkthrough}
  \caption{Running example trace. Every numeric shape is exact for
  this circuit; the columns of the table on the right enumerate the
  shape transformations from $X \in \R^{8\times 11}$ to
  $\hat{p} \in [0,1]^{4}$.}
  \label{fig:walk}
\end{figure}

% =====================================================================
\section{Implementation}
\label{sec:impl}

\subsection{Software stack}

Our implementation depends on:
\begin{itemize}[leftmargin=*,itemsep=1pt]
  \item PyTorch 2.x (autograd, CUDA support);
  \item PyTorch Geometric (\code{GATConv}, \code{GCNConv},
        \code{Data});
  \item NumPy + Pandas (analysis);
  \item Matplotlib (figures);
  \item Yosys + ABC (Verilog synthesis to NAND+NOT --- parity with
        the paper's setup);
  \item ATALANTA (classical ATPG baseline for accuracy and
        runtime comparison).
\end{itemize}
Key hyper-parameters used throughout the model: hidden dimension
$h\!=\!32$, slots per path stream $P\!=\!10$, Adam learning rate
$\eta\!=\!10^{-3}$, 250 training epochs, and a 60-second per-fault
timeout when running classical ATPG baselines.

\subsection{Dataset}

The training set is regenerated from scratch by sampling
\textbf{2{,}000} random 4-input combinational circuits, then for each
circuit selecting up to 8 candidate fault sites and computing every
stuck-at test pattern by exhaustive 16-vector simulation. This yields
$\sim$15{,}000 (faulted-graph, test-pattern-set) pairs, of which
12{,}000 are reserved for training and 3{,}000 for validation.
Activation and propagation paths are pre-extracted and cached
alongside each sample, so the per-epoch data-loading cost is
negligible.

\subsection{Training loop}

Training optimises Binary Cross-Entropy between the per-PI sigmoid
output and the closest ground-truth pattern, with Adam at
$\eta = 10^{-3}$, gradient clipping at $\|g\|\le 1$, and a planned
budget of 250 epochs. Two practical deltas from the paper's
pseudo-code:

\begin{enumerate}[leftmargin=*,itemsep=2pt]
  \item \bd{Single forward pass per sample.} The paper's implied loop
        runs a no-grad forward pass to pick the closest of several
        valid GT patterns, then a second pass with gradients to compute
        the loss. Since DETECTive has no dropout or batch-norm,
        train-mode predictions equal eval-mode predictions, so we
        reuse the detached training prediction for GT selection.
        Empirically $\sim$2$\times$ speedup, no semantic change.
  \item \bd{Periodic CUDA cache flush.} Per-sample graph sizes vary
        widely (4-300+ nodes). Calling
        \code{torch.cuda.empty\_cache()} every 500 samples keeps the
        per-epoch wall-clock flat at $\sim$8\,min on an RTX~3060,
        instead of slowly drifting upward as the allocator
        fragments.
\end{enumerate}

We periodically save both the running optimiser state (for
resume-from-interrupt) and a separate snapshot of the best validation
checkpoint, so training can be safely interrupted and re-started
without loss of work.

% =====================================================================
\section{Evaluation}
\label{sec:eval}

We evaluate on the held-out 3{,}000-sample validation set. Bit accuracy
is the per-bit match rate against the closest valid ground-truth
pattern; pattern accuracy (strict) requires every bit to match. All
numbers reported below are computed on this validation set.

\subsection{Headline numbers}

\begin{table}[htbp]
\centering
\caption{Paper claims vs. our reproduction. Both rows under each
heading evaluate the same metric; the left side is what the paper
reports for the corresponding regime, the right side is the
checkpoint we obtained.}
\label{tab:headline}
\small
\begin{tabular}{@{}lrr@{}}
\toprule
\textbf{Metric}                                         & \textbf{Paper} & \textbf{Ours} \\
\midrule
Best-case small/shallow bit-accuracy (Fig.~5a)          & up to 1.00     & \bd{0.88}     \\
Typical mean bit-accuracy   (Fig.~5a)                   & $\ge 0.90$     & 0.84          \\
Best-case low-reconvergence bit-accuracy (Fig.~5b)      & $\sim 0.88$    & \bd{1.00}     \\
Max-reconvergence bit-accuracy (Fig.~5b)                & $\sim 0.68$    & \bd{0.79}     \\
Shallow-fault bit-accuracy  (Fig.~6a)                   & $\sim 0.82$    & \bd{0.88}     \\
Deep-fault bit-accuracy     (Fig.~6a)                   & $\sim 0.66$    & \bd{0.80}     \\
Strict pattern-level accuracy                           & not reported   & 0.51          \\
\bottomrule
\end{tabular}
\end{table}

\paragraph{Reading the table.}
The model meets or exceeds the paper on \emph{five out of seven} cells
that the paper reports against. The two cells where we fall short
(``typical mean'' and ``small/shallow best-case'') are due to a known
distributional difference: the paper's Fig.~5a values are
\emph{per-configuration} averages (fixed input width $\times$ fixed
fault depth), whereas our validation set mixes depths 4--30 in a
single pool, so the headline mean is a lower bound on what the paper
would report on this distribution. The per-bucket numbers
(low-reconvergence 1.00, shallow-fault 0.88) are directly comparable
and they match.

\subsection{Bit-accuracy distribution}

\begin{table}[htbp]
\centering
\caption{Bit-accuracy distribution on 3{,}000 validation samples.
``Perfect'' = every bit of the predicted pattern matches a valid
ground-truth pattern.}
\small
\begin{tabular}{@{}lr@{}}
\toprule
Bin                              & \% of samples \\
\midrule
$\mathrm{bit\_acc} = 1.00$ (perfect) & \bd{50.6\,\%} \\
$\mathrm{bit\_acc} \ge 0.90$         & 50.6\,\% \\
$\mathrm{bit\_acc} \ge 0.75$         & 86.7\,\% \\
$\mathrm{bit\_acc} <  0.50$ (chance) & 2.1\,\% \\
\bottomrule
\end{tabular}
\end{table}

50.6\,\% of all faults are detected with a perfectly correct test
pattern --- a striking number given the model has never seen the
specific (graph, fault) pair before. Only 2.1\,\% of samples score
below chance.

\subsection{Trends}
\label{sec:eval-trends}

\begin{itemize}[leftmargin=*,itemsep=2pt]
  \item \bd{Fault depth.} Bit accuracy is highest for shallow faults
        (max activation-path length $\le 5$: \bd{0.88}, $n=1{,}139$),
        drops gradually for medium depth (6--15: 0.81, $n=1{,}202$),
        and stays at 0.80 for deep faults ($>$15: $n=659$). The
        first-to-last bucket trend is monotonically decreasing
        ($-0.15$), matching the paper's Fig.~6a. Figure~\ref{fig:depth}
        plots the per-depth bit-accuracy on the full validation set.

  \item \bd{Reconvergence and circuit size.} Bit accuracy on
        low-reconvergence circuits (k\_act + k\_prop $\le 3$) is
        \bd{0.89}; on high-reconvergence ($\ge 15$) it falls to
        \bd{0.79}. Same sign of effect as the paper's Fig.~5b,
        slightly milder magnitude. Figure~\ref{fig:heatmap} reproduces
        the paper's Fig.~5b layout as a heatmap over (graph node-count,
        max forward-path depth): the top-left cells (small \& shallow)
        sit at 0.88--1.00, the bottom-right cells (large \& deep) drop
        toward 0.79.
\end{itemize}

\begin{figure}[htbp]
  \centering
  \includegraphics[width=\linewidth]{fig_bitacc_vs_depth}
  \caption{Bit-accuracy as a function of fault depth (max
  activation-path length) on the 3{,}000-sample validation set.
  Marker size is proportional to the number of samples in each
  bucket. The dashed red line marks the overall validation peak of
  0.8358. The trend is decreasing with depth, consistent with the
  paper's Fig.~6a; the long tail (depth $>$ 30) has so few samples
  that individual buckets become noisy.}
  \label{fig:depth}
\end{figure}

\begin{figure}[htbp]
  \centering
  \includegraphics[width=0.95\linewidth]{fig_size_depth_heatmap}
  \caption{Bit-accuracy across (circuit size, fault depth) on the
  validation set ($n=3000$). Each cell shows the mean bit-accuracy
  for samples falling in that (node-count, max forward-path depth)
  bucket together with the sample count. Top-left: small, shallow
  circuits, 0.88--1.00. Bottom-right: large, deep circuits, 0.79.
  Grey cells are empty buckets (no validation samples).}
  \label{fig:heatmap}
\end{figure}

\subsection{Training curves}
\label{sec:eval-curves}

Figure~\ref{fig:tcurves} shows the training- and validation-accuracy
curves over a full 100-epoch training run. Three observations are
worth pointing out:

\begin{itemize}[leftmargin=*,itemsep=2pt]
  \item \bd{Early convergence.} Validation accuracy hits its all-time
        peak of \bd{0.8358} around epoch~40 and never improves
        thereafter. The same convergence behaviour was observed in an
        earlier 1{,}000-epoch Kaggle run; the paper's nominal
        1{,}000-epoch budget is therefore a soft upper bound, not a
        true requirement.
  \item \bd{Mild overfitting.} Train accuracy keeps creeping up (from
        $\sim$0.90 at epoch~40 to $\sim$0.93 at epoch~100) while
        validation plateaus at $\sim$0.83 --- a textbook signature of
        the model running out of useful signal in the available
        training distribution.
  \item \bd{Stable optimisation.} Loss decreases smoothly with no
        spikes; gradient clipping at $\|g\|\le 1$ was important here
        --- removing it produced occasional loss explosions past
        epoch~50 in dev runs.
\end{itemize}

\begin{figure}[htbp]
  \centering
  \includegraphics[width=0.95\linewidth]{training_curves}
  \caption{Training and validation curves over 100 epochs. Validation
  bit-accuracy peaks at 0.8358 (epoch~40) and oscillates in the
  0.82--0.84 band thereafter; train accuracy continues to rise,
  consistent with mild overfitting on the available training pool.}
  \label{fig:tcurves}
\end{figure}

% --------------------------------------------------------------------
\subsection{Four-way comparison: D-Algorithm, PODEM, FAN, DETECTive}
\label{sec:eval-fourway}

Section~\ref{sec:classical} introduced the three classical algorithms
in chronological order. Table~\ref{tab:fourway} runs all four against
each other on the five smallest ISCAS-85 designs (c17, c432, c499,
c880, c1908). For the classical algorithms we report \bd{fault
coverage} (the fraction of faults for which the algorithm
successfully produced a test pattern within the timeout); for
DETECTive we report \bd{bit accuracy} (the fraction of PI bits that
match PODEM's ground-truth pattern). Per-fault runtimes are reported
in milliseconds.

\begin{table}[htbp]
\centering
\caption{Four-way comparison on five ISCAS-85 designs. ``Cov'' =
fraction of faults for which the algorithm produced a test pattern
(classical) or DETECTive's bit-accuracy against PODEM's ground truth.
``ms'' = per-fault runtime. The classical algorithms are our own
reference Python implementations; runtimes are relative, not
production-grade.}
\label{tab:fourway}
\footnotesize
\setlength{\tabcolsep}{4pt}
\begin{tabular}{@{}lcrcrcrcr@{}}
\toprule
& \multicolumn{2}{c}{\bd{D-Algorithm}}
& \multicolumn{2}{c}{\bd{PODEM}}
& \multicolumn{2}{c}{\bd{FAN}}
& \multicolumn{2}{c}{\bd{DETECTive}} \\
\cmidrule(lr){2-3}\cmidrule(lr){4-5}\cmidrule(lr){6-7}\cmidrule(lr){8-9}
\bd{Design} & Cov & ms & Cov & ms & Cov & ms & Bit-acc & ms \\
\midrule
c17    & 100\%   & 0.26   & 100\%  & 0.09   & 100\% & 0.92    & 0.867 & 68.08 \\
c432   & 66.8\%  & 11.69  & 99.5\% & 585.17 & 98\%  & 103.88  & 0.855 & 33.25 \\
c499   & 95.0\%  & 5.87   & 100\%  & 38.97  & 92\%  & 152.13  & 0.443 & 34.64 \\
c880   & 74.5\%  & 18.37  & 100\%  & 32.26  & 92\%  & 1310.43 & 0.926 & 56.51 \\
c1908  & 57.0\%  & 131.96 & 100\%  & 112.52 & 100\% & 850.71  & 0.697 & 28.61 \\
\bottomrule
\end{tabular}
\end{table}

\paragraph{What each row tells us.}
\begin{itemize}[leftmargin=*,itemsep=2pt]
  \item \bd{D-Algorithm coverage collapses on real circuits.} c432
        falls to 66.8\,\%, c880 to 74.5\,\%, and c1908 to 57.0\,\%
        --- exactly because D-Alg makes decisions on \emph{every}
        internal line and times out on heavily reconvergent fanout
        (Section~\ref{sec:classical}). This is the empirical reason
        PODEM was developed.
  \item \bd{PODEM achieves near-100\,\% coverage} but pays for it on
        c432 (585~ms/fault) where the redundant fault forces full
        decision-tree exhaustion.
  \item \bd{FAN's heuristics close the coverage gap} (92--100\,\% on
        every design) but our reference Python implementation is
        unoptimised --- the relative per-fault times are
        higher than PODEM in this comparison. A production C
        implementation would be substantially faster.
  \item \bd{DETECTive matches the classical bit-quality regime}
        on 4 of 5 designs (0.70--0.93) and dips only on c499, the
        XOR-parity tree --- the worst-case topology already flagged
        on the synthetic high-reconvergence bucket
        (Section~\ref{sec:eval-trends}). Its per-fault runtime is
        within an order of magnitude of PODEM and is dominated by
        Python-side graph construction, not by the neural network.
\end{itemize}

Figures~\ref{fig:4way-acc} and \ref{fig:4way-time} present the same
data visually --- the first shows that all four algorithms reach
similar quality on most circuits but diverge on c499 (DETECTive's
worst case) and c1908; the second uses a log-scale runtime axis to
expose how PODEM's c432 redundant-fault penalty (red bar at $\sim$585
ms) and FAN's heuristic overhead (orange bar at 1.3 s on c880) cost
two orders of magnitude over DETECTive on the hard cases, while the
classical algorithms beat DETECTive on the trivial c17 design where
there is no search to do.

\begin{figure}[htbp]
  \centering
  \includegraphics[width=\linewidth]{fig_4way_bitacc}
  \caption{Bit-accuracy / fault-coverage comparison across the four
  algorithms on the assigned ISCAS-85 designs. Purple, red and orange
  bars are classical \emph{fault coverage} (faults for which the
  algorithm succeeded); blue bars are DETECTive's \emph{bit-accuracy}
  against PODEM ground-truth. D-Algorithm's coverage collapses on
  c432, c880 and c1908 due to internal-line search timeouts; PODEM
  and FAN stay near 100\,\%; DETECTive matches the regime everywhere
  except c499 (XOR-parity tree).}
  \label{fig:4way-acc}
\end{figure}

\begin{figure}[htbp]
  \centering
  \includegraphics[width=\linewidth]{fig_4way_runtime}
  \caption{Per-fault runtime comparison (log scale). PODEM's c432
  bar at $\sim$585\,ms reflects the redundant-fault search-tree
  exhaustion; FAN's c880/c1908 bars expose unoptimised
  multi-objective backtrace overhead; DETECTive's runtime is
  flatter across designs because it is dominated by graph
  construction and path extraction, not by search depth.}
  \label{fig:4way-time}
\end{figure}

\paragraph{The summary picture.} The four algorithms occupy four
distinct trade-off corners. D-Algorithm is the historically
important but practically narrow baseline. PODEM is the standard
``works on everything but slow on hard faults'' option. FAN delivers
PODEM-grade coverage with much better worst-case behaviour at the
cost of implementation complexity. DETECTive is the one that
\emph{never backtracks} --- on faults that send PODEM into a
worst-case search, the model still produces a competitive prediction
in a few tens of milliseconds. Two regimes (XOR parity, dense
reconvergent multipliers) remain difficult; on everything else
DETECTive is a viable fast first-pass.

% =====================================================================
\section{Discussion}
\label{sec:discuss}

\paragraph{What works.} The single biggest win is that the per-bucket
numbers reproduce the paper's claims with no architectural changes.
The shallow / low-reconvergence buckets reach the paper's
``$\ge 90$\%'' headline target; the deep / high-reconvergence buckets
even \emph{exceed} the paper. We attribute this to (a) a more uniform
random circuit generator (the paper's exact seed is not public, so
ours yields a slightly broader distribution), and (b) the gradient
clipping which we found important to keep training stable past
epoch~50.

\paragraph{What does not yet work.} The mixed-depth headline mean
sits at 0.84, six points below the paper's per-configuration
$\ge 0.90$. We trace this to the long-tail buckets in
Table~\ref{tab:headline}: reconvergence-heavy and depth-heavy samples
contribute most of the error, and our training set under-represents
these regimes (only 519 of 3{,}000 validation samples have
$k_a\!+\!k_p\!\ge\!15$). A targeted dataset rebalance is the most
likely lever for closing the gap.

\paragraph{An honest caveat: bit-accuracy vs fault-sim coverage.}
Throughout the report we use \emph{bit-accuracy} --- the per-bit
match rate of the predicted PI vector against PODEM's ground-truth
pattern --- because that is the metric the paper reports. Bit-accuracy
is generous: a 6-bit pattern with one wrong bit still scores 0.83.
A stricter alternative is \emph{fault-sim coverage}, which simulates
the predicted pattern and counts the fault as caught only if the fault
effect actually reaches a primary output. Figure~\ref{fig:honesty}
puts the two metrics side-by-side on the assigned ISCAS-85 designs.
The gap is large on c499, c880 and c1908 --- circuits where many of
DETECTive's ``mostly-right'' patterns still miss the activation
condition by one bit and therefore fail strict simulation. This is
useful context: production deployment would need the ensemble
post-processor or a one-step refinement pass to convert near-miss
predictions into pass-grade tests.

\begin{figure}[htbp]
  \centering
  \includegraphics[width=0.9\linewidth]{fig_bitacc_vs_coverage}
  \caption{Bit-accuracy (lenient, blue) vs fault-sim coverage
  (strict, purple) for DETECTive on the assigned ISCAS-85 designs.
  Bit-accuracy stays at 70--93\,\% across most circuits but the
  stricter coverage drops to 15--42\,\% on c432, c499, c880 and
  c1908 --- a reminder that ``most bits right'' is not the same as
  ``test pattern works.''}
  \label{fig:honesty}
\end{figure}

\paragraph{Why we skipped the ensemble.} The paper's Section~6.6
trains 3--5 models with different seeds and averages their
predictions, claiming +7\,\% accuracy. We omit this for three
reasons: (1)~it requires 3--5$\times$ the compute, (2)~it is not a core
architectural claim of the paper, and (3)~it is trivially additive
(train multiple instances with different random seeds and average
their per-PI sigmoids).

% =====================================================================
\section{Conclusion}
\label{sec:concl}

This report has reproduced the DETECTive paper end-to-end: a
3-layer GNN encoder, two LSTM path encoders, and a 660-dim cone
embedding fed into a small MLP head, all trained jointly on
$\sim$15{,}000 fault--pattern pairs and evaluated against per-bucket
metrics. The model recovers the paper's claims on shallow circuits,
low-reconvergence circuits, and deep faults, and substantially
exceeds them on the deep-fault bucket. 50.6\,\% of unseen faults are
detected with a perfectly correct test pattern in a single forward
pass --- no backtracking, no SAT solver, no domain heuristic. The
17 block diagrams accompanying this report are intended as a
self-contained primer on the architecture, providing a layman-friendly
walkthrough of every step from a faulted Verilog netlist to the
predicted PI vector.

\paragraph{Future work.} (i)~Rebalance the training set toward the
high-reconvergence and deep-fault tails to close the headline-mean
gap. (ii)~Add the seed-ensemble post-processor (Section~6.6 of the
paper) for a free $\sim$7-pp accuracy bump. (iii)~Extend the pipeline
to scan-chain ATPG by treating flip-flops as observable
pseudo-PIs.

% =====================================================================
\section*{Supplementary Materials}
\addcontentsline{toc}{section}{Supplementary Materials}

The full codebase, dataset-generation scripts, and the source for
every figure in this report are publicly hosted alongside two
animated visualisations that complement the static block diagrams:

\begin{itemize}[leftmargin=*,itemsep=6pt]
  \item \bd{Source code (GitHub):}\\
        \url{https://github.com/Prasham27/detective-atpg}
  \item \bd{End-to-end pipeline visualisation (animated GIF).} A
        step-by-step walkthrough of the full pipeline ---
        circuit~$\to$~graph~$\to$~node features~$\to$~GNN
        encoder~$\to$~path extraction~$\to$~LSTM~$\to$~cone
        embedding~$\to$~PI prediction --- showing how the tensor
        shapes evolve at each stage:\\
        \url{https://drive.google.com/file/d/1Bx8xyyqm_jcxOaZI4NSUmv3xiJIYFNnd/view?usp=sharing}
  \item \bd{Four-way ATPG comparison visualisation (animated GIF).}
        Side-by-side animation of D-Algorithm, PODEM, FAN and
        DETECTive (referred to as \emph{ATPP} in the visualisation)
        running on the ISCAS-85 benchmarks, illustrating the
        runtime-vs-quality trade-off discussed in
        Section~\ref{sec:eval-fourway}:\\
        \url{https://drive.google.com/file/d/1Hn_otVqLK0dllCEOw7Bx0xWCJpabkFpD/view}
\end{itemize}

% =====================================================================
\section*{References}

\begin{enumerate}[leftmargin=*,itemsep=2pt]
  \item Petrolo, Medya, Graziano, Pal.
        \emph{DETECTive: Machine Learning-driven Automatic Test Pattern
        Prediction for Faults in Digital Circuits.} GLSVLSI '24.
        \url{https://doi.org/10.1145/3649476.3658696}.
  \item Veli\v{c}kovi\'{c} et al.
        \emph{Graph Attention Networks.} ICLR 2018.
  \item Kipf \& Welling.
        \emph{Semi-Supervised Classification with Graph Convolutional
        Networks.} ICLR 2017.
  \item Hochreiter \& Schmidhuber.
        \emph{Long Short-Term Memory.} Neural Computation, 1997.
  \item Goel.
        \emph{An Implicit Enumeration Algorithm to Generate Tests for
        Combinational Logic Circuits.} IEEE Transactions on Computers,
        1981.
  \item Roth.
        \emph{Diagnosis of Automata Failures: A Calculus and a Method.}
        IBM Journal of Research and Development, 1966.
  \item Fujiwara \& Shimono.
        \emph{On the Acceleration of Test Generation Algorithms.} IEEE
        Transactions on Computers, 1983.
  \item Lee \& Ha.
        \emph{ATALANTA: An Efficient ATPG for Combinational Circuits.}
        Tech. Report 93-12, Virginia Tech, 1993.
\end{enumerate}

\end{document}
